Fix \(P\in\mathcal M\).  We first construct a full-data lift and then compare the two decision problems.  Put \(Z=(X,A)\).  Since \(\mathcal X\) is finite, both \(\mathcal X\times[0,1]\) and \([0,1]\) are standard Borel spaces.  Kallenberg's regular-conditional-distribution theorem, in the exact form recorded by \(\texttt{lem:standard-borel-regular-conditional-law}\), supplies a Borel probability kernel \(K\) from \(\mathcal X\times[0,1]\) to \([0,1]\) that is a conditional law of \(Y\) given \(Z\).  Write \(K_x(a,C)=K((x,a),C)\).  By \(\texttt{ass:conditional-density-law}\), for Borel \(B,C\subseteq[0,1]\),
\[
 P\{X=x,A\in B,Y\in C\}
 =p_x\int_B K_x(a,C)\pi_x(a)\,\mathrm da. \tag{1}
\]
The map
\[
 m_x^0(a)=\int_{[0,1]}y\,K_x(a,\mathrm dy) \tag{2}
\]
is Borel: kernel integration is Borel first for indicators, then for nonnegative simple functions, and hence for the bounded identity by monotone approximation.  It is a version of \(\mathbb E_P[Y\mid A=a,X=x]\).  By \(\texttt{ass:holder-regression}\), \(\mu_x^P\) is a continuous version of the same conditional mean.  Thus
\[
 E_x=\{a\in[0,\bar\delta]:m_x^0(a)\ne\mu_x^P(a)\} \tag{3}
\]
is Borel.  Since \(p_x>0\) by \(\texttt{ass:stratum-mass}\), uniqueness of conditional expectations together with \(\texttt{ass:conditional-density-law}\) yields
\[
 \int_{E_x}\pi_x(a)\,\mathrm da=0. \tag{4}
\]
Modify the kernel only on this exposure-null set by
\[
 K_x^\star(a,C)=
 \mathbf 1\{a\notin E_x\}K_x(a,C)
 +\mathbf 1\{a\in E_x\}\mathbf 1\{\mu_x^P(a)\in C\}. \tag{5}
\]
For each Borel \(C\), (5) is Borel in \(a\), and for each \(a\) it is a probability measure in \(C\).  Equations (3)--(5) give
\[
 \int y\,K_x^\star(a,\mathrm dy)=\mu_x^P(a)
 \quad(0\leq a\leq\bar\delta), \tag{6}
\]
whereas (4) gives
\[
 p_x\int_BK_x^\star(a,C)\pi_x(a)\,\mathrm da
 =p_x\int_BK_x(a,C)\pi_x(a)\,\mathrm da. \tag{7}
\]

We realize the modified kernel by an explicit measurable quantile.  Let \(D=\mathbb Q\cap[0,1]\), including \(1\), set \(F_x(a,t)=K_x^\star(a,[0,t])\), and define
\[
 g_x(a,u)=\inf\{r\in D:F_x(a,r)\geq u\}. \tag{8}
\]
The set in (8) is nonempty because \(F_x(a,1)=1\).  For every real \(t\),
\[
 \{(a,u):g_x(a,u)<t\}
 =\bigcup_{\substack{r\in D\\r<t}}
 \{(a,u):u\leq F_x(a,r)\}, \tag{9}
\]
a countable union of Borel sets, so \(g_x\) is jointly Borel.  Continuity from above of the probability measure \(K_x^\star(a,\cdot)\) makes \(t\mapsto F_x(a,t)\) right-continuous.  Consequently
\[
 \lambda\{u:g_x(a,u)\leq t\}=F_x(a,t), \tag{10}
\]
where \(\lambda\) is Lebesgue measure on \([0,1]\).  Indeed, the two inclusions follow respectively from rationals decreasing to \(t\) and, when \(u>F_x(a,t)\), from right continuity at \(t\).  Since the intervals \([0,t]\) generate \(\mathcal B([0,1])\), (10) says that \(u\mapsto g_x(a,u)\) pushes \(\lambda\) forward to \(K_x^\star(a,\cdot)\).

On a probability space carrying \((X,A)\) with its \(P\)-margin and an independent \(U\sim\operatorname{Unif}[0,1]\), define simultaneously
\[
 Y(a)=g_X(a,U)\quad(0\leq a\leq1),
 \qquad Y=g_X(A,U). \tag{11}
\]
Finiteness of \(\mathcal X\) and joint Borel measurability of the \(g_x\)'s make \((a,\omega)\mapsto Y(a,\omega)\) jointly measurable.  Equation (11) verifies \(\texttt{ass:consistency}\) on one probability-one event, and independence gives \(A\perp U\mid X\), namely \(\texttt{ass:exchangeability}\).  By (10), the conditional law of \(Y\) given \((X,A)=(x,a)\) is \(K_x^\star(a,\cdot)\).  Equations (1) and (7), first on measurable rectangles and then on the generated product sigma-field, therefore show that the observed \((X,A,Y)\)-law is exactly \(P\).  Finally, (6) and (10) imply
\[
 m_x^{\mathrm F}(a)
 =\int_0^1g_x(a,u)\,\mathrm du
 =\int y\,K_x^\star(a,\mathrm dy)
 =\mu_x^P(a),
 \qquad 0\leq a\leq\bar\delta. \tag{12}
\]
The right side is continuous by \(\texttt{ass:holder-regression}\), so \(\texttt{ass:full-data-response-continuity}\) holds.  Thus (11), with the member-specific carrier \(\mathcal U=[0,1]\), is a member of \(\mathcal M^{\mathrm F}\) by \(\texttt{def:full-data-model-class}\).  This proves that the observed-margin map \(\Gamma\) satisfies
\[
 \Gamma(\mathcal M^{\mathrm F})=\mathcal M. \tag{13}
\]

By \(\texttt{prop:causal-bridge}\), for every member and every declared threshold,
\[
 \psi_{\delta_n}(P^{\mathrm F})
 =\theta_{\delta_n}\{\Gamma(P^{\mathrm F})\}. \tag{14}
\]
If \(T\) is any estimator measurable with respect to \(O_1,\ldots,O_n\), its law under \((P^{\mathrm F})^{\otimes n}\) is its law under \(\Gamma(P^{\mathrm F})^{\otimes n}\).  Equations (13)--(14) therefore imply
\[
 \sup_{P^{\mathrm F}\in\mathcal M^{\mathrm F}}
 \mathbb E_{(P^{\mathrm F})^{\otimes n}}
 |T-\psi_{\delta_n}(P^{\mathrm F})|
 =\sup_{P\in\mathcal M}
 \mathbb E_{P^{\otimes n}}|T-\theta_{\delta_n}(P)|. \tag{15}
\]
The estimator classes are identical, because both consist of measurable maps of the observed sample.  Taking their infima in (15), using \(\texttt{def:causal-frontier-criteria}\) and the definition of \(R_n^\star\), gives
\[
 R_{n,\mathrm F}^\star=R_n^\star. \tag{16}
\]
For any observed-sample interval \(C_n\), the same argument gives
\[
 \inf_{P^{\mathrm F}\in\mathcal M^{\mathrm F}}
 (P^{\mathrm F})^{\otimes n}\{\psi_{\delta_n}(P^{\mathrm F})\in C_n\}
 =\inf_{P\in\mathcal M}P^{\otimes n}\{\theta_{\delta_n}(P)\in C_n\} \tag{17}
\]
and
\[
 \sup_{P^{\mathrm F}\in\mathcal M^{\mathrm F}}
 \mathbb E_{(P^{\mathrm F})^{\otimes n}}[\operatorname{len}(C_n)]
 =\sup_{P\in\mathcal M}
 \mathbb E_{P^{\otimes n}}[\operatorname{len}(C_n)]. \tag{18}
\]
Hence the coverage constraints select the same intervals and their objectives coincide.  Taking the infimum in (18) over the common feasible set in (17) proves
\[
 L_{n,\mathrm F}^\star=L_n^\star. \tag{19}
\]
Equations (13), (16), and (19) are the asserted conclusions.